% not printed: handwritten call number "QA153 W35 v.1"; library stamps "MATH.-STAT. LIBRARY" and "M777593"
\begin{center}
{\Large \gesperrt{PREFACE.}}\\[1ex]
\rule{4em}{0.4pt}
\end{center}

In view of the development that algebra has undergone in recent decades, a comprehensive presentation and linking of the various theoretical considerations and manifold applications may well be useful, even after the textbook of \gesperrt{Serret}, which was excellent for its time.

For years I have cherished the plan of such an undertaking, which indeed appeared vast and far-reaching and required a good deal of preparatory work. Only after I had several times traversed the field as a whole in university lectures, and had treated individual parts in greater detail, did I resolve to set about the execution of the work, of which the first volume now lies before the reader in completed form.

It was my intention to provide a textbook which, without presupposing much prior knowledge, would introduce the reader to modern algebra and also lead him on to the higher and more difficult parts, in which interest in the subject first becomes truly alive. At the same time, the necessary tools, the elementary as well as the higher ones, were to be derived from the course of the development itself, in order to make the presentation as independent of other textbooks as possible.

There are two things that have become of quite special importance for the most recent development of algebra; these are, on the one hand, group theory, which is gaining ever greater dominance and whose ordering and clarifying influence can be felt everywhere, and then the intervention of number theory.
